% Part: sets-functions-relations
% Chapter: sets
% Section: unions-and-intersections

\documentclass[../../../include/open-logic-section]{subfiles}

\begin{document}

\olfileid[id]{sfr}{set}{uni}
\olsection{Gabungan dan Irisan}

\begin{explain}
Dalam \olref[sfr][set][bas]{sec}, kita memperkenalkan definisi
himpunan melalui abstraksi, yaitu definisi berbentuk
$\Setabs{x}{\phi(x)}$. Di sini kita menggunakan suatu sifat~$\phi$,
dan sifat tersebut dapat menyebut himpunan yang sudah kita definisikan.
Sebagai contoh, jika $A$ dan~$B$ adalah himpunan, maka himpunan
$\Setabs{x}{x \in A \lor x \in B}$ terdiri atas semua objek yang
merupakan !!{element}s dari $A$ atau~$B$. Dengan kata lain, himpunan
itu menggabungkan !!{element}s dari $A$ dan~$B$. Kita dapat
memvisualisasikannya seperti dalam \olref{fig:union}; daerah yang
diarsir menunjukkan !!{element}s kedua himpunan $A$ dan~$B$ secara
bersama-sama.

\begin{figure}
  \olasset{assets/diagrams/union.tikz}
  \caption{Gabungan $A \cup B$ dari dua himpunan adalah himpunan
   !!{element}s dari $A$ bersama dengan anggota dari~$B$.}
  \ollabel{fig:union}
\end{figure}

Operasi pada himpunan ini---yakni menggabungkan keduanya---sangat
berguna dan umum, sehingga kita memberinya nama dan lambang formal.
\end{explain}

\begin{defn}[Gabungan]
\emph{Gabungan} dua himpunan $A$ dan $B$, yang ditulis $A \cup B$,
adalah himpunan semua objek yang merupakan !!{element}s dari $A$,
dari $B$, atau dari keduanya.
\[
A \cup B = \Setabs{x}{x \in A \lor x \in B}
\]
\end{defn}

\begin{ex}
Karena banyaknya pencantuman !!{element}s tidak berpengaruh, gabungan
dua himpunan yang mempunyai !!a{element} bersama hanya memuat
!!{element} tersebut satu kali. Sebagai contoh,
$\{ a, b, c\} \cup \{ a, 0, 1\} = \{a, b, c, 0, 1\}$.

Gabungan sebuah himpunan dengan salah satu himpunan bagiannya sama
dengan himpunan yang lebih besar: $\{a, b, c \} \cup \{a \} =
\{a, b, c\}$.

Gabungan sebuah himpunan dengan himpunan kosong identik dengan
himpunan itu: $\{a, b, c \} \cup \emptyset = \{a, b, c \}$.
\end{ex}

\begin{prob}
Buktikan bahwa jika $A \subseteq B$, maka $A \cup B = B$.
\end{prob}

\begin{explain}
Kita juga dapat meninjau operasi yang ``dual'' terhadap gabungan.
Operasi ini membentuk himpunan semua !!{element}s yang merupakan
!!{element}s dari~$A$ sekaligus !!{element}s dari~$B$. Operasi ini
disebut \emph{irisan} dan dapat digambarkan seperti dalam
\olref{fig:intersection}.
\begin{figure}
  \olasset{assets/diagrams/intersection.tikz}
  \caption{Irisan $A \cap B$ dari dua himpunan adalah himpunan
    !!{element}s yang dimiliki keduanya secara bersama.}
  \ollabel{fig:intersection}
\end{figure}
\end{explain}

\begin{defn}[Irisan]
\emph{Irisan} dua himpunan $A$ dan $B$, yang ditulis $A \cap B$,
adalah himpunan semua objek yang merupakan !!{element}s dari $A$ dan
juga dari~$B$.
\[
A \cap B = \Setabs{x}{x \in A \land x \in B}
\]
Dua himpunan disebut \emph{saling asing} jika irisannya kosong. Ini
berarti keduanya tidak mempunyai !!{element}s bersama.
\end{defn}

\begin{ex}
Jika dua himpunan tidak mempunyai !!{element}s bersama, irisannya
kosong: $\{ a, b, c\} \cap \{ 0, 1\} = \emptyset$.

Jika dua himpunan mempunyai !!{element}s bersama, irisannya adalah
himpunan semua anggota tersebut: $\{a, b, c \} \cap \{a, b, d \} =
\{a, b\}$.

Irisan sebuah himpunan dengan salah satu himpunan bagiannya sama
dengan himpunan yang lebih kecil: $\{a, b, c\} \cap \{a, b\} =
\{a, b\}$.

Irisan setiap himpunan dengan himpunan kosong adalah kosong:
$\{a, b, c \} \cap \emptyset = \emptyset$.
\end{ex}

\begin{prob}
Buktikan secara ketat bahwa jika $A \subseteq B$, maka $A \cap B = A$.
\end{prob}

\begin{explain}
Kita juga dapat membentuk gabungan atau irisan lebih dari dua himpunan.
Cara yang anggun untuk menanganinya secara umum adalah sebagai berikut:
misalkan semua himpunan yang ingin kita gabungkan (atau irisankan)
dihimpun ke dalam satu himpunan. Kita kemudian dapat mendefinisikan
gabungan semua himpunan semula sebagai himpunan semua objek yang
termasuk dalam paling sedikit satu !!{element} dari himpunan tersebut,
dan mendefinisikan irisannya sebagai himpunan semua objek yang termasuk
dalam setiap !!{element} dari himpunan tersebut.
\end{explain}

\begin{defn}
Jika $A$ adalah himpunan dari himpunan-himpunan, maka $\bigcup A$
adalah himpunan !!{element}s dari !!{element}s dari~$A$:
\begin{align*}
\bigcup A & = \Setabs{x}{x \text{ termasuk dalam !!a{element} dari } A},
\text{ yaitu,}\\
& = \Setabs{x}{\text{terdapat } B \in A
  \text{ sedemikian sehingga } x \in B}
\end{align*}
\end{defn}

\begin{defn}
Jika $A$ adalah himpunan dari himpunan-himpunan, maka $\bigcap A$
adalah himpunan objek yang dimiliki bersama oleh semua anggota
dari~$A$:
\begin{align*}
\bigcap A & = \Setabs{x}{x \text{ termasuk dalam setiap !!{element} dari } A},
\text{ yaitu,}\\
 & = \Setabs{x}{\text{untuk setiap } B \in A, x \in B}
\end{align*}
\end{defn}

\begin{ex}
Misalkan $A = \{ \{ a, b \}, \{ a, d, e \}, \{ a, d \} \}$.
Maka $\bigcup A = \{ a, b, d, e \}$ dan $\bigcap A = \{ a \}$.
\end{ex}
\begin{prob}
	Tunjukkan bahwa jika $A$ adalah himpunan dan $A \in B$, maka
	$A \subseteq \bigcup B$.
\end{prob}

Kita juga dapat melakukan hal yang sama untuk barisan himpunan $A_1$,
$A_2$, \dots
\begin{align*}
\bigcup_i A_i & = \Setabs{x}{x \text{ termasuk dalam salah satu } A_i}\\
\bigcap_i A_i & = \Setabs{x}{x \text{ termasuk dalam setiap } A_i}.
\end{align*}

Jika kita mempunyai suatu \emph{himpunan indeks} bagi himpunan-himpunan
itu, yaitu himpunan $I$ sedemikian sehingga kita meninjau $A_i$ untuk
setiap $i \in I$, kita juga dapat menggunakan singkatan berikut:
\begin{align*}
	\bigcup_{i \in I} A_i & = \bigcup \Setabs{A_i }{i \in I}\\
	\bigcap_{i \in I} A_i & = \bigcap\Setabs{A_i}{i \in I}
\end{align*}

Terakhir, kita mungkin ingin meninjau himpunan semua !!{element}s
dalam~$A$ yang tidak terdapat dalam~$B$. Hal ini dapat digambarkan
seperti dalam \olref{difference}.

\begin{figure}
  \olasset{assets/diagrams/difference.tikz}
  \caption{Selisih $A \setminus B$ dari dua himpunan adalah himpunan
    !!{element}s dari~$A$ yang bukan merupakan !!{element}s dari~$B$.}
  \ollabel{difference}
\end{figure}

\begin{defn}[Selisih]
\emph{Selisih himpunan}~$A \setminus B$ adalah himpunan semua
!!{element}s dari $A$ yang bukan merupakan !!{element}s dari~$B$,
yaitu
\[
A\setminus B = \Setabs{x}{x\in A \text{ dan } x \notin B}.
\]
\end{defn}

\begin{prob}
	Buktikan bahwa jika $A \subsetneq B$, maka
	$B \setminus A \neq \emptyset$.
\end{prob}

\end{document}
